\section{Hamiltonian Flights}
\label{sec:cses-hamiltonian-flights}

\problemstatement{Given a directed graph of $n$ cities ($n\le20$) and $m$
flights, count the number of routes from city $1$ to city $n$ that visit
\emph{every} city exactly once (Hamiltonian paths). Output the count modulo
$10^9+7$.}

\begin{patternbox}
Hamiltonian path counting is NP-hard in general, but $n\le20$ invites
\textbf{bitmask DP}: the state is (set of visited cities, current city), and
$2^n\cdot n$ states are affordable.
\end{patternbox}

\subsection*{DP over visited subsets}

Let $\textit{dp}[\textit{mask}][v]$ be the number of routes that start at
city $1$, have visited exactly the set $\textit{mask}$ (which must include
$1$ and $v$), and currently sit at $v$. Base case:
$\textit{dp}[\{1\}][1]=1$. Transition: from state $(\textit{mask},u)$, for
each flight $u\to v$ with $v\notin\textit{mask}$, add
$\textit{dp}[\textit{mask}][u]$ to
$\textit{dp}[\textit{mask}\cup\{v\}][v]$. The answer is
$\textit{dp}[\text{full}][n]$ -- all cities visited, ending at $n$.

\begin{ideabox}[The Visited Set Is the State]
A Hamiltonian path cannot revisit a city, so the exact set already visited
(plus where you are) fully determines valid continuations. Encoding that set
as a bitmask makes the DP finite: $2^n$ subsets times $n$ endpoints.
\end{ideabox}

\subsection*{Algorithm}

\begin{enumerate}
  \item $\textit{dp}[1\ll0][1]=1$ (city $1$ is bit $0$).
  \item For each mask (increasing) and each city $u$ in it with
        $\textit{dp}>0$: extend along flights $u\to v$ with $v$ unvisited.
  \item Output $\textit{dp}[(1\ll n)-1][n-1]$.
\end{enumerate}

\subsection*{Dry run}

\dryrun{$n=3$; flights $1\!\to\!2$, $2\!\to\!3$, $1\!\to\!3$.}

\begin{itemize}
  \item Only $1\to2\to3$ visits all three ending at $3$; $1\to3$ skips
        city $2$.
  \item Count $=1$.
\end{itemize}

\subsection*{C++ implementation}

\begin{lstlisting}
#include <bits/stdc++.h>
using namespace std;
const long long MOD = 1e9 + 7;

int main() {
    int n, m;
    cin >> n >> m;
    vector<vector<int>> adj(n);                    // 0-indexed cities
    for (int i = 0; i < m; i++) {
        int a, b; cin >> a >> b;
        adj[a - 1].push_back(b - 1);
    }

    int full = 1 << n;
    vector<vector<long long>> dp(full, vector<long long>(n, 0));
    dp[1][0] = 1;                                  // visited {city 0}, at city 0
    for (int mask = 1; mask < full; mask++)
        for (int u = 0; u < n; u++) {
            if (!(mask & (1 << u)) || dp[mask][u] == 0) continue;
            for (int v : adj[u])
                if (!(mask & (1 << v)))
                    dp[mask | (1 << v)][v] =
                        (dp[mask | (1 << v)][v] + dp[mask][u]) % MOD;
        }
    cout << dp[full - 1][n - 1] << "\n";
    return 0;
}
\end{lstlisting}

\begin{complexitybox}
\textbf{Time Complexity.} $O(2^n\cdot n^2)$ in the worst case ($n$ endpoints,
$n$ neighbours). \textbf{Space Complexity.} $O(2^n\cdot n)$ for the DP table.
\end{complexitybox}

\begin{takeawaybox}
Hamiltonian counting is bitmask DP with state (visited set, current node);
small $n$ is what makes the exponential state space tractable. This
subset-DP template also solves the travelling salesman problem and
assignment-style problems.
\end{takeawaybox}
